\documentclass[openany,fontset=fandol]{ctexbook}
\usepackage{amsmath,amssymb,amsthm}
\usepackage[unicode,hidelinks]{hyperref}
\setlength{\emergencystretch}{2em}
\raggedbottom
\newcommand{\OO}{\mathcal O}
\newcommand{\FF}{\mathcal F}
\newcommand{\GG}{\mathcal G}
\newcommand{\PP}{\mathbf P}
\DeclareMathOperator{\Spec}{Spec}
\DeclareMathOperator{\coprof}{coprof}
\newcommand{\tagitem}[2]{\par\medskip\noindent\textbf{#1\quad #2}\enspace}
\begin{document}
\title{代数几何基础第五卷预备笔记\\法文覆盖审计：泛截面的局部性质与光滑性}
\author{第89批\quad 法、英、中文三方校订}
\date{2026年11月23日}
\maketitle

\section*{范围与接续}
本批逐页对照法文物理10--12页及英编稿物理9--10页，从命题2.8到推论2.13；
第12页后半的定理2.14起留待下一批。既有合订稿已有对应译文，本批提供校订文本，
不把这几页重复计成新译页，也不修改1158页基线。
沿用第88批记号：$P=\PP^n_k$，$P^\vee$ 为超平面参数空间，$\eta$ 是其泛点；
$Y=X\times_P H$ 为关联族，$h:Y_\eta\to X$ 为泛超平面截面的自然投影。
正式出版 EGA IV.1 仍另行接续。

\section*{20.2\quad 泛超平面截面：局部性质（续）}
\tagitem{命题}{2.8}
设 $\FF$ 是 $X$ 上的凝聚模层，$y\in Y_\eta$，$x=h(y)$。
对诺特局部环 $A$ 上的有限型模 $M$，取性质 $\mathcal P(M)$ 为下列任意一种：
\begin{enumerate}
 \item[(i)] $\coprof M\le n$（参考待补）；
 \item[(ii)] $M$ 满足 $(S_k)$（参考待补）；
 \item[(iii)] $M$ 是 Cohen--Macaulay 模；
 \item[(iv)] $M$ 约化（参考待补）；
 \item[(v)] $M$ 整（参考待补）。
\end{enumerate}
则 $\GG_{\eta,y}$ 满足 $\mathcal P$，当且仅当 $\FF_x$ 满足 $\mathcal P$。
\begin{proof}
依据原稿所引第6段关于正则局部同态的结论，并用 $Y_\eta\to X$ 是正则态射，
因而 $\OO_{X,x}\to\OO_{Y_\eta,y}$ 是正则局部同态。结合（2.3），可作下一推论。
\end{proof}
\noindent\emph{译注。} 法文第10页明确写 $Y_\eta\to X$，英编稿写成 $Y_\eta\to Y$。
后一个态射也正则，因为它是由 $\Spec k(\eta)\to P^\vee$ 基变换所得的局部化；
但要推得所写的局部同态，须再复合射影丛 $Y\to X$。
原稿的第6段引用及四处“参考待补”仍保留；英稿疑此处是否指 EGA IV，
不能把编者疑问当成作者给出的确定卷次。性质 $\mathcal P$ 与射影空间 $P$ 分写。

\tagitem{推论}{2.9}
沿用（2.8）的记号。令 $Z$ 为 $\FF_x$ 不满足 $\mathcal P$ 的点集。
在该坏点集可构造时，$\GG_\eta$ 处处满足 $\mathcal P$，当且仅当 $f(Z)$ 有限，
等价于 $\dim\overline{f(Z)}\le 0$。
\begin{proof}
由（2.8），$h^{-1}(Z)$ 正是 $\GG_\eta$ 不满足 $\mathcal P$ 的点集。
对可构造的 $Z$ 用（2.3），得它为空当且仅当 $f(Z)$ 有限。
作者在此注明 $h$ 指 $Y_\eta\to X$，并指出（2.8）中原用字母 $P$ 指两个对象。
\end{proof}
\noindent\emph{译注。} 法、英原稿均写成维数等于零，空集也应计入，故改为不大于零。
（2.3）在这里需要坏点集的可构造性，正文显式写出所用条件；
不能把“任何由若干闭点组成的集合”都当成有限集。

\tagitem{推论}{2.10}
使 $Y_\eta$ 约化、分别使它局部为整概形的条件。%
\footnote{法10和英9均只有这一行提纲。下面划去的说明不能当成这个推论的完整陈述或证明。}

\tagitem{推论}{2.11}
设 $y\in Y_\eta$，$x=h(y)$。
$Y_\eta$ 在 $y$ 处正则、满足 $(R_k)$ 或正规，分别当且仅当 $X$ 在 $x$ 处满足同一性质。
对应地，令 $Z$ 为 $X$ 中的非正则点、不满足 $(R_k)$ 的点或非正规点所成的集合；
$Y_\eta$ 处处满足相应性质，分别等价于 $f(Z)$ 有限，即 $\dim\overline{f(Z)}\le0$。
\begin{proof}
沿用（2.8）和（2.9）的论证。原稿提示，须给出各个 $Z$ 为闭集的参考，
因为用（2.3）至少需要它可构造。
\end{proof}
\noindent\emph{译校补记。}
在这里的有限型域上情形，正则轨迹与正规轨迹为开；
而局部环不满足 $(R_k)$，等价于存在高度不超过 $k$ 的非正则素点向该点特化。
若这样一个非正则素点存在，则包含它的某个非正则轨迹不可约分支的泛点
高度也不超过 $k$。反之，这样的泛点向该点特化便违反 $(R_k)$。
因此相应坏点集是非正则轨迹中这些分支的并，仍为闭。%
\footnote{法11页的手写改动清楚保留 $(R_k)$，英编稿有一次作 $E_k$；空集的修正同（2.9）。所述闭性说明是译校补充，不冒充原稿已列资料。}

作者另说：（2.10）和（2.11）只讨论绝对性质，不涉及相应的几何性质。
下面考察几何性质，也可以趁此另起一节。%
\footnote{法10末已有一段类似文字，被划去并以“及下述2.11”做了增注；法11在（2.11）证明后有未划的完整表述。英编稿亦将完整段放在（2.11）后，并加小标题。划删稿的独立信息仅是举“几何约化”“几何整”为例，末句“下面我们将……”未完。}

\section*{几何性质（英编稿标题）}
\tagitem{定理}{2.12}
设 $f:X\to P$ 非分歧。令 $y\in Y_\eta$，$x=h(y)$。
则 $X$ 在 $x$ 处对 $k$ 光滑，当且仅当 $Y_\eta$ 在 $y$ 处对 $k(\eta)$ 光滑。%
\footnote{法11页原打字为“浸入”，被划去，手写改为“非分歧”。英编稿的陈述采用修订后的假设，但正文证明没有同步改写。非分歧不推出 Zariski 局部浸入，如非平凡有限可分域扩张的谱到点。}

\begin{proof}[订正后的证明]
先证从截面光滑推出 $X$ 光滑。取 $k$ 的代数闭包 $\bar k$，
将构造换到 $\bar k$，并记新对偶空间泛点为 $\bar\eta$。
其函数域是 $\bar k(P^\vee)$（即换底后的对偶射影空间的函数域），
是 $k(\eta)$ 的扩域；新泛截面是原泛截面沿此扩域的基变换。
取 $y$ 上方一点 $\bar y$，设它在 $X_{\bar k}$ 的像为 $\bar x$。
由光滑性随域扩张保持，$\bar y$ 处的局部环正则。
泛截面向 $X_{\bar k}$ 的投影平坦（射影丛复合取泛纤维的局部化），
故由平坦局部同态上的正则性下降，$X_{\bar k}$ 在 $\bar x$ 处正则。
$\bar k$ 完美，有限型 $\bar k$-概形在正则点光滑；再沿域扩张下降光滑性，
得到 $X$ 在 $x$ 处光滑。这一方向不使用非分歧。

反向在 $x$ 的光滑开邻域内进行，并缩小到纯相对维数 $d$ 的一块。
写射影坐标给出的截面为 $s_0,\ldots,s_n$，它们生成 $L=f^*\OO_P(1)$。
在 $f(x)$ 落入的坐标片内，取 $s_0$ 平凡化 $L$，令 $T_i=s_i/s_0$。
这样每个参数 $[a_0:\cdots:a_n]$ 给出方程
\[
 a_0+\sum_{i=1}^n a_i T_i=0.
\]
因为 $f$ 非分歧，$\Omega^1_{P/k}$ 向 $\Omega^1_{X/k}$ 的拉回映射满射，
所以 $dT_i$ 在这一邻域生成秩 $d$ 的余切模。
对每个点，要求既满足上述方程又有 $\sum a_i\,dT_i=0$，
便对 $n+1$ 个参数系数施加 $d+1$ 个独立线性条件：
微分映射的秩为 $d$，而 $a_0$ 只在值条件中以系数一出现。
故坏参数连同其上的点，在此邻域构成秩 $n-d$ 的核丛的射影化。
若秩为零，该射影化为空；否则总维数是
\[
 d+(n-d-1)=n-1<\dim P^\vee .
\]
因此它在参数空间的像闭包不含 $\eta$。
在泛超平面的零点，微分便不为零；在光滑底空间内用一条方程的雅可比判据，
可知该泛截面在 $y$ 处光滑。
\end{proof}

\noindent\emph{原证明异文。}
法11--12页及英10页沿用了浸入情形的坐标法：
在 $P$ 的开集 $U$ 中用 $f_1,\ldots,f_p$ 切出 $X$，
再附加 $\sum_iS_iT_i-1$，欲证明
\[
 df_1,\ldots,df_p,\sum_iS_i\,dT_i
\]
线性无关。原稿写商模
$\Omega^1_{U/k}/\sum_{i=1}^n\OO_Udf_i$ 非零并列条件 $p\le n$；
这里求和应至 $p$，在微分独立处商秩为 $n-p$，非零须 $p<n$。
稿上进一步声称：任一非零局部自由模的一组生成元 $a_i$，
其泛线性组合 $\sum_i S_i a_i$ 在整个扩域概形上都不消失。
这个辅助断言不成立：平凡线丛由 $1,t$ 生成，
$S_1+S_2t$ 却在 $t=-S_1/S_2$ 处为零。
原稿最后另说 $p=n$ 时（该局部零维情形）泛超平面截面为空。
订正证明只在截面的零点考察微分，把值和一阶微分条件同时计入，不使用这条错误断言。%
\footnote{现代参考：\href{https://stacks.math.columbia.edu/tag/05AW}{Stacks Project 05AW}（平坦态射下几何正则性下降），
\href{https://stacks.math.columbia.edu/tag/0FD4}{0FD4}（Bertini 与关联族方法），
\href{https://stacks.math.columbia.edu/tag/07R6}{07R6}（正则态射的定义与性质）。
这里的维数证明明确写出非分歧态射情形，不把任何仅对闭浸入陈述的结论直接套用。}

\tagitem{推论}{2.13}
仍假设 $f$ 非分歧。设 $Z$ 为 $X$ 不在 $k$ 上光滑的点集。
则 $Y_\eta$ 对 $k(\eta)$ 光滑，当且仅当 $Z$ 有限。
特别是 $X$ 光滑时，$Y_\eta$ 也光滑。
\begin{proof}
光滑轨迹开，所以 $Z$ 闭。由（2.12）、（2.3），所求等价于 $f(Z)$ 有限。
非分歧的有限型态射拟有限，每个纤维的点集有限，故这里 $f(Z)$ 有限恰等价于 $Z$ 有限。
\end{proof}
\noindent 下一处为同一页后半定理2.14（法物理12，英物理10）。
法文在该定理列表后手写补了局部环性质的双向关系，不能仅据两种底本的换页位置断成两条。

\end{document}
